\documentclass[11pt]{beamer}
\usepackage{multimedia}
\usepackage[ngerman]{babel}
\beamertemplatenavigationsymbolsempty
\useoutertheme[subsection=false]{miniframes}
%\usefonttheme{professionalfonts} % using non standard fonts for beamer
\setbeamertemplate{footline}[frame number]{}
\setbeamersize{text margin left=-5pt,text margin right=-5pt} 

\usepackage[T1]{fontenc}

\usepackage{mathtools}
\usepackage{amsfonts, amsmath, amssymb, amsthm}
\usepackage{pifont}
\usepackage{ wasysym }
\usepackage{MnSymbol}
\usepackage{xspace}
\usepackage{ stmaryrd }
\usepackage{bbold} %for mathbb 1
\usepackage{chngcntr}

\usepackage{graphicx}

\usepackage{pgf, tikz}
\usetikzlibrary{automata, arrows, arrows.meta, positioning, decorations.pathreplacing, decorations.pathmorphing, calc, fit}
\usepackage{varwidth}
\usepackage{color}
\usepackage{geometry}
\usepackage{tcolorbox}
\tcbuselibrary{xparse}
\usepackage{multirow}

\setbeamersize{text margin left=12.5mm,text margin right=10mm}
%---------------------------------------------------------------
%colors

%base colors
\definecolor{wblau}{HTML}{B1D9EE}
\definecolor{blau}{HTML}{569BC0}
\definecolor{sblau}{HTML}{1E6B94}
\definecolor{worange}{HTML}{FFDFBA}
\definecolor{orange}{HTML}{FFBB6B}
\definecolor{sorange}{HTML}{E98F25}
\definecolor{gelb}{HTML}{c44848}

\definecolor{wgelb}{HTML}{F9E79F}

\definecolor{theorie}{HTML}{db6b94}
\definecolor{praxis}{HTML}{FFBB6B}
\definecolor{zukunft}{HTML}{569BC0}
%highlighting
\definecolor{wgruen}{HTML}{CAEA9C}
\definecolor{gruen}{HTML}{8FC542}
\definecolor{sgruen}{HTML}{2E4E00}
\definecolor{wviolett}{HTML}{D68FAF}
\definecolor{violett}{HTML}{B43C73}
\definecolor{sviolett}{HTML}{470021}

\definecolor{rot}{rgb}{.8,0,0}
\definecolor{rot2}{HTML}{E27D60}

\definecolor{turkis}{rgb}{0,.7,.7}

\definecolor{grau}{rgb}{0.47,0.47,0.47}
\definecolor{grau2}{rgb}{0.2,0.2,0.2}

\definecolor{titlebg}{HTML}{213845}
\definecolor{titlefg}{HTML}{F7E1C8}

%---------------------------------------------------------------------------

\usetheme{Singapore}
%	AnnArbor | Antibes | Bergen |
%	Berkeley | Berlin | Boadilla |
%	boxes | CambridgeUS | Copenhagen |
%	Darmstadt | default | Dresden |
%	Frankfurt | Goettingen |Hannover |
%	Ilmenau | JuanLesPins | Luebeck |
%	Madrid | Malmoe | Marburg |
%	Montpellier | PaloAlto | Pittsburgh |
%	Rochester | Singapore | Szeged |
%	Warsaw
\usecolortheme{beaver}
%	albatross | beaver | beetle |
%	crane | default | dolphin |
%	dove | fly | lily | orchid |
%	rose |seagull | seahorse |
%	sidebartab | structure |
%	whale | wolverine
%\usefonttheme{professionalfonts}
%	default | professionalfonts | serif |
%	structurebold | structureitalicserif |
%	structuresmallcapsserif
%\useinnertheme{rectangles}
%	circles | default | inmargin |
%	rectangles | rounded
%\useoutertheme{infolines}
%	default | infolines | miniframes |
%	shadow | sidebar | smoothbars |
%	smoothtree | split | tree
\makeatletter
\AtBeginDocument{%
  {
    \usebeamercolor{section in head/foot}
  }
  \pgfdeclareverticalshading{beamer@headfade}{\paperwidth}
  {%
    color(0cm)=(white);
    color(1.25cm)=(white)%
  }
  \setbeamercolor{section in head/foot}{bg=}
}
\addtoheadtemplate{\pgfuseshading{beamer@headfade}\vskip-1.25cm}{}
\makeatother

\makeatletter
\setbeamertemplate{frametitle}{
    \ifbeamercolorempty[bg]{frametitle}{}{\nointerlineskip}%
    \@tempdima=\textwidth%
    \advance\@tempdima by\beamer@leftmargin%
    \advance\@tempdima by\beamer@rightmargin%
    \vspace*{-13pt} %%%%%%%%%%%%% For example insert shift to right
    % \hspace*{-8.1pt} %%%%%%%%%%%%% For example insert shift to right
    \begin{beamercolorbox}[sep=.435cm,center,wd=\the\@tempdima]{frametitle}
        \usebeamerfont{frametitle}%
        \vbox{}\vskip-1ex%
        \if@tempswa\else\csname beamer@ftecenter\endcsname\fi%
        \strut\insertframetitle\strut\par%
        {%
            \ifx\insertframesubtitle\@empty%
            \else%
            {\usebeamerfont{framesubtitle}\usebeamercolor[fg]{framesubtitle}\insertframesubtitle\strut\par}%
            \fi
        }%
        \vskip-1ex%
        \if@tempswa\else\vskip-.3cm\fi% set inside beamercolorbox... evil here...
    \end{beamercolorbox}%
}
\makeatother
\setbeamercolor{frametitle}{fg=titlefg,bg=titlebg}
\setbeamercolor{title}{fg=titlefg,bg=titlebg}
%------------------------------------------------------------------
%boxes
\newtcolorbox{obox}[1][]{
colframe=worange,
colback=worange,
left=5pt,
right=5pt,
#1
}
\newtcolorbox{bbox}[1][]{
colframe=wblau,
colback=wblau,
left=5pt,
right=5pt,
#1
}

\DeclareTColorBox{tbox}{ O{} m }{
colframe=white,
colback=white,
grow to left by=0.8cm,
grow to right by=0.15cm,
leftupper=.8cm,
title=#2,
coltitle=sorange,
#1
}
\DeclareTColorBox{dbox}{ O{} m }{
colframe=white,
colback=white,
grow to left by=0.8cm,
grow to right by=0.15cm,
leftupper=.8cm,
title=#2,
coltitle=sblau,
#1
}

\newcommand*\titem{%
  \item[\color{orange}{\textbullet}]}
\newcommand*\ditem{%
  \item[\color{blau}{\textbullet}]}
  \newcommand*\nitem{%
  \item[\color{titlebg}{\textbullet}]}


%---------------------------------------------------------------------------
\pgfdeclarelayer{TransLayer}
\pgfsetlayers{TransLayer,main}

\tikzset{
  AutomatonState/.style = {
      inner sep = 0pt,
      line width = 4pt,
      draw = titlebg!75!white,
      minimum size = 20pt,
      initial text = {Start},
      initial where = left,
      accepting/.style = {
        accepting by arrow,
        accepting text = {Ende},
        accepting where = right,
    },
  },
  AutomatonEdge/.style={
      draw=titlebg,
      -{Stealth[inset=0pt, width=5pt, length=5pt]},
      >={Stealth[inset=0pt, width=5pt,length=5pt]},
      shorten >=0pt,
      shorten <=-.5pt,
      line width=2pt,
  },
  AutomatonLabel/.style={ text=black, scale=1.2},
  Automaton/.style = {
      node distance=22mm,
      on grid,
      >=Latex,
      initial text=,
      initial where=below left,
      every node/.style={AutomatonLabel},
      every state/.style = {AutomatonState},
      every edge/.style = {AutomatonEdge},
  },
}
%---------------------------------------------------------------------------
\newcommand{\ca}{A^*_c}
\newcommand{\A}{\mathcal{A}}
\newcommand{\ma}{A^*}
\newcommand{\cont}{\mathop{c}}
\newcommand{\decont}{\mathop{\reflectbox{c}}}

\newcommand{\ini}{\mathop{in}}
\newcommand{\wt}{\mathop{wt}}
\newcommand{\out}{\mathop{out}}

\newcommand{\ldl}{\text{LDL}}
\newcommand{\wldl}{\ifmmode\text{wLDL}\else$\text{wLDL}$\fi\xspace}
\newcommand{\wldlp}{\ifmmode{\text{wLDL}}^{\mkern-3mu+}\else${\text{wLDL}}^{\mkern-3mu+}$\fi}
\newcommand{\wafa}{\ifmmode \text{WAFA} \else WAFA \fi}
\newcommand{\mso}{\ifmmode \text{MSO} \else MSO \fi}
\newcommand{\lvl}{\ifmmode \text{LvlA}\else LvlA\fi\xspace}

\newcommand{\ser}[1]{S\langle#1\rangle}
\newcommand{\chr}[1]{\mathbb{1}_{#1}}
\newcommand{\bhv}[1]{\llbracket #1 \rrbracket}
\newcommand{\ms}[1]{\{\mkern-7mu\{#1\}\mkern-7mu\}}
\newcommand{\sem}[1]{\scalebox{.85}{|}\mkern-2mu\scalebox{.85}{|}#1\scalebox{.85}{|}\mkern-2mu\scalebox{.85}{|}}
\newcommand{\tsem}[1]{\scalebox{.85}{(}\mkern-4mu\scalebox{.85}{(}#1\scalebox{.85}{)}\mkern-4mu\scalebox{.85}{)}}
\newcommand{\asem}[1]{\{\mkern-4mu|#1|\mkern-4mu\}}
\newcommand{\csem}[2]{\sem{#2}_{#1}}

\newcommand{\fend}{\text{end}}
\newcommand{\fedge}{\text{E}}
\newcommand{\ftrue}{\text{true}}
\newcommand{\ffalse}{\text{false}}
\newcommand{\co}{%
  \mathbin{%
    \vcenter{\offinterlineskip
      \vspace{4.6pt}%
      \ialign{##\cr$\cdot$\cr\noalign{\kern-1pt}$\mkern-.6mu,$\cr}%
    }%
  }%
}

\newcommand{\colorcircle}[1]{%
  \raisebox{0.25ex}{\tikz{\draw[line width=2pt, color=#1] (0,0) circle (2pt);}}%
}

\newcommand{\colorcross}[1]{%
  \raisebox{0.25ex}{\tikz{
    \draw[line width=2pt, color=#1] (-2pt,-2pt) -- (2pt,2pt);
    \draw[line width=2pt, color=#1] (-2pt,2pt) -- (2pt,-2pt);
  }}%
}
\newcommand{\yes}{\colorcircle{gruen}}
\newcommand{\no}{\colorcross{rot2}}
\newcommand{\comp}{\mathbin{\ooalign{$\hidewidth ;\hidewidth$\cr$\phantom{+}$}}}
\newcommand{\colordot}[1]{\tikz{\fill[opacity=0.8, #1] (0,0) circle (2pt);}}
\newcommand{\gr}[1]{\textcolor{gruen}{#1}}
\newcommand{\vi}[1]{\textcolor{violett}{#1}}
\newcommand{\bl}[1]{\textcolor{blau}{#1}}
\newcommand{\bra}[1]{\raisebox{-1pt}{\scalebox{1.1}{\textless}}\mkern+2mu #1\mkern+2mu \raisebox{-1pt}{\scalebox{1.1}{\textgreater}}\mkern+5mu}
\newcommand{\ubra}[1]{\raisebox{-1pt}{\scalebox{1.1}{[}}\mkern+2mu #1\mkern+2mu \raisebox{-1pt}{\scalebox{1.1}{]}}\mkern+5mu}
\mathtoolsset{centercolon=true}

\newcommand\blfootnote[1]{%
  \begingroup
  \renewcommand\thefootnote{}\footnote{#1}%
  \addtocounter{footnote}{-1}%
  \endgroup
}
\newcommand{\high}[3]{\only<-\number\numexpr#1-1\relax>{#3}\only<#1->{\textcolor{#2}{#3}}}
\newcommand{\shigh}[4]{\only<-\number\numexpr#1-1\relax,\number\numexpr#2+1\relax->{#4}\only<#1-#2>{\textcolor{#3}{#4}}}
\newcommand{\light}[2]{\shigh{#1}{#1}{gruen}{#2}}
 

%---------------------------------------------------------------------------
\title{Nutzen der Automatentheorie zur Verbesserung der Sicherheit durch Verifikation in IT-Systemen von KRITIS-Betreibern}
\author{Gustav Grabolle}
 \date{%
    Halle\\
    August 4, 2025
}
%---------------------------------------------------------------------------

\begin{document}

\maketitle

\begin{frame}
	\begin{tcolorbox}[ colframe=theorie, colback=theorie!35!white, height=2.3cm, valign = center, ]
		\huge
		\begin{center}
			Theoretische Grundlagen
		\end{center}
	\end{tcolorbox}

	\vfill

	\begin{tcolorbox}[ colframe=praxis, colback=praxis!35!white, height=2.3cm, valign = center, ]
		\huge
		\begin{center}
			Praktische Einsicht
		\end{center}
	\end{tcolorbox}


	\vfill

	\begin{tcolorbox}[ colframe=zukunft, colback=zukunft!35!white, height=2.3cm, valign = center, ]
		\huge
		\begin{center}
			KRITIS Einsatz
		\end{center}
	\end{tcolorbox}
\end{frame}

\setbeamercolor{background canvas}{bg=theorie!35!white}
\setbeamercolor{frametitle}{fg=theorie!35!white,bg=titlebg}

\begin{frame}{Was sind diese Automaten?}
	\visible<15->{
	\begin{minipage}{\textwidth}
		\begin{tcolorbox}[colframe=theorie, colback=white, valign=center, halign=center]
			\visible<15->{$\text{\light{15}{A}} \mathrel{\hat{=}} \text{Aufnahme}$} \hfill
			\visible<16->{$\text{\light{16}{B}} \mathrel{\hat{=}} \text{Bearbeitung}$} \hfill
			\visible<17->{$\text{\light{17}{E}} \mathrel{\hat{=}} \text{Einreichung}$}
		\end{tcolorbox}
		}
	\end{minipage}

	\vspace{10pt}

	\begin{minipage}{.48\textwidth}
		\begin{tcolorbox}[colframe=theorie, colback=white, valign=center, halign=center, height=170pt]
			\hspace*{-13pt}
			\begin{tikzpicture}[Automaton,]
				\node[state,initial, accepting] (q1) at (0,0) {};
				\node[state] (q2) [below of = q1] {};

				\draw[->] (q1) edge[bend left] node[AutomatonLabel, right]{\light{15}{A}} (q2);
				\draw[->] (q2) edge[loop below] node[AutomatonLabel, below] {\light{16}{B}} (q2);
				\draw[->] (q2) edge[bend left] node[AutomatonLabel, left]{\light{17}{E}} (q1);

				\only<3,6,9>{\draw[fill=gruen, draw=none] (q1) circle (7pt);}
				\only<4,5,10>{\draw[fill=gruen, draw=none] (q2) circle (7pt);}
				\only<11>{\draw[fill=rot2, draw=none] (q2) circle (7pt);}
				\only<4,10>{\draw[->] (q1) edge[bend left, draw = gruen] node[AutomatonLabel, right]{A} (q2);}
				\only<5>{\draw[->] (q2) edge[loop below, draw = gruen] node[AutomatonLabel, below] {B} (q2);}
				\only<6>{\draw[->] (q2) edge[bend left, draw = gruen] node[AutomatonLabel, left]{E} (q1);}
			\end{tikzpicture}
		\end{tcolorbox}
	\end{minipage}%
	\begin{minipage}{ .04\textwidth }%
		\hfill%
	\end{minipage}%
	\begin{minipage}{.48\textwidth}%
		\visible<2->{
			\begin{tcolorbox}[colframe=theorie, colback=white, valign=center, halign=center, height=60pt]
				\begin{tabular}{lc}
					\visible<2->{\light{4}{A}\light{5}{B}\light{6}{E}}  & \visible<7->{\yes}  \\
					\visible<8->{\light{10}{A}\shigh{11}{11}{rot2}{A}E} & \visible<12->{\no}  \\
					\visible<13->{ABEAE}                                & \visible<14->{\yes} \\
				\end{tabular}
			\end{tcolorbox}
		}
		\visible<19->{
			\begin{tcolorbox}[colframe=theorie, colback=white, height=103pt, valign=center, boxsep=0pt]
				\scriptsize
				\visible<19->{
					\textit{%
						Jede Aufnahme benötigt eine Einreichung, bevor eine neue Aufnahme erfolgt.
					}
				}
				\vfill
				\visible<20->{\[|u|_A = |u|_E \text{ oder } |u|_A = |u|_E + 1\]}
			\end{tcolorbox}
		}
	\end{minipage}
\end{frame}

\begin{frame}{Verifikation mit Modelchecking}
	\vfill
	\begin{center}
		\begin{tikzpicture}[
				node distance=0.5cm and 1.5cm,
				box/.style={
						draw=theorie,
						fill=white,
						line width = 1.5pt,
						rounded corners=3pt,
						minimum width=4cm,
						minimum height=0.8cm,
						align=center,
						inner sep=3pt,
						font=\normalsize,
					},
				arrow/.style={->, thick, draw=theorie}
			]


			% Linke Spalte
			\visible<1->{\node[box] (spec1) {\textbf{Spezifikation}};}
			\visible<4->{\node[box, below=of spec1] (spec2) {Logik};}
			\visible<5->{\node[box, below=of spec2] (spec3) {Automat $A_L$};}

			% Rechte Spalte
			\visible<2->{\node[box, right=of spec1] (real1) {\textbf{Realisation}};}
			\visible<6->{\node[box, below=of real1] (real2) {Programm};}
			\visible<7->{\node[box, below=of real2] (real3) {Automat $A_P$};}

			% Gleichheitszeichen zwischen spec1 und real1
			\visible<3->{\node at ($(spec1)!0.5!(real1)$) {\light{11}{\huge$\mathrel{\hat{=}}$}};}

			% Pfeile
			\visible<4->{\draw[arrow, AutomatonEdge] (spec1) -- (spec2);}
			\visible<5->{\draw[arrow, AutomatonEdge] (spec2) -- (spec3);}
			\visible<6->{\draw[arrow, AutomatonEdge] (real1) -- (real2);}
			\visible<7->{\draw[arrow, AutomatonEdge] (real2) -- (real3);}

			% Schnittfrage-Box
			\visible<8->{
				\node[fit=(spec3.east)(real3.west), inner sep=0pt] (fitbox) {};
				\path let \p1 = (real3.east),
				\p2 = (spec3.west),
				\n1 = {\x1 - \x2} in
				node[box, below=1.2cm of fitbox, minimum height=60pt, minimum width=\n1, text width=\n1-30pt] (check) {
						\only<9->{
							\scriptsize
							\textit{Gibt es einen erfolgreichen Lauf des Programmautomaten,\\ der kein erfolgreicher Lauf des Logikautomaten ist?}\\[10pt]
							\normalsize
							\only<-9>{
								\phantom{
									$L(A_P) \cap \overline{L(A_L)} = \emptyset$?
								}
							}
						}
						\only<10->{
							\normalsize
							$L(A_P) \cap \overline{L(A_L)} = \emptyset$\light{11}{?}
						}
					};
			}

			% Pfeile zur Schnittfrage
			\visible<8->{\draw[arrow, AutomatonEdge] (spec3.south) -- (check.north -| spec3.south);}
			\visible<8->{\draw[arrow, AutomatonEdge] (real3.south) -- (check.north -| real3.south);}
		\end{tikzpicture}
	\end{center}
	\vfill
\end{frame}

\setbeamercolor{background canvas}{bg=praxis!35!white}
\setbeamercolor{frametitle}{fg=praxis!35!white,bg=titlebg}

\begin{frame}{Modelchecking in Aktion}
	\begin{tcolorbox}[
			colframe = praxis,
			colback=white,
			valign=center,
			halign = center,
			left=-18pt,
			right=-18pt,
			top=0pt,
			bottom=0pt,
			boxsep=0pt,
			left skip = 20pt,
			right skip = 20pt,
			% sharp corners,
			arc = 0mm,
			boxrule = 4pt,
		]
		\movie[width=.877\textwidth]
		{\includegraphics[width=.9\textwidth]{../sources/transparent.png}}{../sources/recording.mkv}
	\end{tcolorbox}

\end{frame}


\setbeamercolor{background canvas}{bg=zukunft!35!white}
\setbeamercolor{frametitle}{fg=zukunft!35!white,bg=titlebg}

\begin{frame}{Modelchecking im Finanzwesen}

\end{frame}


\begin{frame}{Modelchecking in Bahnsteuerungssystemen}

\end{frame}

\begin{frame}{Kosten der Softwareverifikation}
	\vfill
	\begin{center}
		\begin{tikzpicture}[
				node distance=0.8cm and 1cm,
				theorie/.style={blue!70!black},
				box/.style={
						draw=theorie,
						fill=white,
						line width=1.5pt,
						rounded corners=3pt,
						minimum width=9cm,
						minimum height=2.4cm,
						align=center,
						inner sep=3pt,
						font=\normalsize,
						text centered,
						anchor=north west,
					},
				subbox/.style={
						draw=theorie,
						fill=white,
						line width=1.5pt,
						inner sep=0pt,
						minimum height=0.8cm,
						text centered,
						font=\normalsize,
						anchor=north west,
					},
				doublearrow/.style={
						<->, thick, theorie,
					},
			]

			% Überschrift oben
			\node[font=\bfseries] (softwarelabel) at (0,3.2) {Software};

			% Oberes Rechteck
			\node[box] (softwarebox) at (0,0) {};

			% Breite des äußeren Rechtecks und Breite der drei inneren Teile
			\def\outerwidth{9}
			\def\outerheight{2.4}
			\def\innerwidth{\outerwidth/3}
			\def\innerheight{\outerheight}

			% Drei innere Teile oben (Programm 1, 2, 3)
			\foreach \i/\text in {0/Programm 1,1/Programm 2,2/Programm 3} {
			\node[subbox, minimum width=\innerwidth cm, minimum height=\innerheight cm]
			(at={($(softwarebox.north west) + (\i*\innerwidth cm,0)$)}) {\textbf{\text}};
			}

			% Vertikale Linien zur Unterteilung (parallel und durchgehend)
			\draw[theorie, line width=1.5pt]
			($(softwarebox.north west) + (3cm,0)$) -- ($(softwarebox.south west) + (3cm,0)$);
			\draw[theorie, line width=1.5pt]
			($(softwarebox.north west) + (6cm,0)$) -- ($(softwarebox.south west) + (6cm,0)$);

			% Unteres Rechteck (Verifikation)
			\node[box, anchor=north west] (verifikationsbox) at (0,-3.5) {};

			% Drei innere Teile unten (Zertifikat 1, 2, 3)
			\foreach \i/\text in {0/Zertifikat 1,1/Zertifikat 2,2/Zertifikat 3} {
			\node[subbox, minimum width=\innerwidth cm, minimum height=\innerheight cm]
			(at={($(verifikationsbox.north west) + (\i*\innerwidth cm,0)$)}) {\textbf{\text}};
			}

			% Vertikale Linien zur Unterteilung unten
			\draw[theorie, line width=1.5pt]
			($(verifikationsbox.north west) + (3cm,0)$) -- ($(verifikationsbox.south west) + (3cm,0)$);
			\draw[theorie, line width=1.5pt]
			($(verifikationsbox.north west) + (6cm,0)$) -- ($(verifikationsbox.south west) + (6cm,0)$);

			% Überschrift unten
			\node[font=\bfseries] (verificationlabel) at (4.5,-6.2) {Verifikation};

			% Doppelpfeile zwischen korrespondierenden Teilen
			\foreach \x in {1.5,4.5,7.5} {
					\draw[doublearrow]
					($(softwarebox.south west) + (\x cm,0)$) --
					($(verifikationsbox.north west) + (\x cm,0)$);
				}

		\end{tikzpicture}
	\end{center}
	\vfill
\end{frame}
\begin{frame}{Der Preis der Sicherheit}

\end{frame}

\begin{frame}{Who will check the checkmen?}

\end{frame}
\end{document}
